% Source: physical PDF pages 491--495; printed pages 468--472.
% Appendix A begins after Section 23.8 on physical p. 491 and ends before
% the Appendix B heading on physical p. 495.
% Figures: none. Source uncertainty resolved at 600 dpi: Eq. (23.A.21)
% explicitly uses q(0)=q and q(T)=q', while Eqs. (23.A.17), (23.A.22),
% and (23.A.23) use the interval -T/2 to T/2.
\chapterappendix{A}{Euclidean Path Integrals}
\phantomsection\label{app:23.A}

This appendix will outline the use of Euclidean path integrals in quantum field theory. As mentioned in \hyperref[sec:9.1]{Section 9.1}, it is possible to formulate quantum field theory in a four-dimensional Euclidean spacetime. Instead of going into the non-trivial analytic continuation needed to calculate S-matrix elements in this approach, here we shall illustrate the use of Euclidean path integrals by addressing a problem for which they are naturally suited.

We consider a set of Hermitian canonical variables $Q_a$ and $P_a$, with commutation relations
\begin{equation}
 [Q_a,P_b]=\ii\delta_{ab},
 \tag{23.A.1}\label{eq:23.A.1}
\end{equation}
\begin{equation}
 [Q_a,Q_b]=[P_a,P_b]=0.
 \tag{23.A.2}\label{eq:23.A.2}
\end{equation}
In quantum field theory the index $a$ is understood, as in \hyperref[sec:9.1]{Section 9.1}, to consist of a spatial position $\mathbf x$ and any discrete Lorentz and species indices $m$, and the Kronecker delta in Eq.~\eqref{eq:23.A.1} is understood as $\delta_{\mathbf x m,\mathbf y n}=\delta^3(\mathbf x-\mathbf y)\delta_{mn}$. We define eigenstates of the $Q_a$
\begin{equation}
 Q_a\ket{q}=q_a\ket{q},
 \tag{23.A.3}\label{eq:23.A.3}
\end{equation}
normalized so that
\begin{equation}
 \braket{q'}{q}=\delta(q'-q)\equiv\prod_a\delta(q'_a-q_a),
 \tag{23.A.4}\label{eq:23.A.4}
\end{equation}
and likewise for eigenstates $\ket{p}$ of the $P_a$.

The problem we consider is the calculation of the matrix element
\begin{equation}
 F(q',q;T)\equiv\bra{q'}\exp\bigl(-H[Q,P]T\bigr)\ket{q},
 \tag{23.A.5}\label{eq:23.A.5}
\end{equation}
where $H$ is the Hamiltonian, and $T$ is an arbitrary positive constant. One application is to the study of ground-state energies. If the smallest eigenvalue of $H$ is $E_0$, with eigenvector $\ket{0}$, then for $T\to\infty$
\[
 F(q',q;T)\to\braket{q'}{0}\braket{0}{q}\exp(-E_0T),
\]
so
\begin{equation}
 E_0=-\lim_{T\to\infty}\left(\frac{\ln F(q',q;T)}{T}\right).
 \tag{23.A.6}\label{eq:23.A.6}
\end{equation}
Also, we can calculate the partition function of statistical mechanics from the trace:
\begin{equation}
 Z(\beta)\equiv\operatorname{Tr}\exp(-\beta H)
 =\int\left[\prod_a \dd q_a\right]F(q,q;\beta),
 \tag{23.A.7}\label{eq:23.A.7}
\end{equation}
where $1/\beta$ is the temperature.

To derive a path integral formula for $F(q',q;T)$ we define Euclidean time dependent operators
\begin{equation}
 Q_a(t)\equiv e^{Ht}Q_ae^{-Ht},
 \qquad
 P_a(t)\equiv e^{Ht}P_ae^{-Ht},
 \tag{23.A.8}\label{eq:23.A.8}
\end{equation}
and corresponding left- and right-eigenstates
\begin{equation}
 \ket{q,t}\equiv\exp(Ht)\ket{q},
 \qquad
 \bra{q,t}\equiv\bra{q}\exp(-Ht),
 \tag{23.A.9}\label{eq:23.A.9}
\end{equation}
\begin{equation}
 \ket{p,t}\equiv\exp(Ht)\ket{p},
 \qquad
 \bra{p,t}\equiv\bra{p}\exp(-Ht),
 \tag{23.A.10}\label{eq:23.A.10}
\end{equation}
such that
\begin{equation}
 Q_a(t)\ket{q,t}=q_a\ket{q,t},
 \qquad
 \bra{q,t}Q_a(t)=q_a\bra{q,t},
 \tag{23.A.11}\label{eq:23.A.11}
\end{equation}
and
\begin{equation}
 P_a(t)\ket{p,t}=p_a\ket{p,t},
 \qquad
 \bra{p,t}P_a(t)=p_a\bra{p,t}.
 \tag{23.A.12}\label{eq:23.A.12}
\end{equation}
One difference between this and the usual Minkowskian formalism is that the `time' evolution of the operators is governed by a non-unitary similarity transformation \eqref{eq:23.A.8}, so there is no simple relation between the right-eigenstates $\bra{q,t}$ of $Q(t)$ and the Hermitian adjoint of the left-eigenstates $\ket{q,t}$, except at $t=0$.

In this language, the definition \eqref{eq:23.A.5} of $F(q',q;T)$ may be rewritten
\begin{equation}
 F(q',q;T)=\braket{q',T/2}{q,-T/2}.
 \tag{23.A.13}\label{eq:23.A.13}
\end{equation}
Let us first calculate $\braket{q',t+\dd t}{q,t}$ for infinitesimal $\dd t$. Adopting the convention that $H(Q,P)$ is written with all $Q$s to the left of all $P$s, we have
\[
 \braket{q',t+\dd t}{q,t}
 =\bra{q',t}\exp\bigl(-H(q',P)\dd t\bigr)\ket{q,t}.
\]
\begingroup
\emergencystretch=1em
Let us expand $\ket{q,t}$ in a complete set of eigenstates of the $P_a(t)$ operators. From Eq.~\eqref{eq:23.A.1} we have as usual
\par
\endgroup
\[
 \braket{q,t}{p,t}
 =\prod_a\frac{\exp(\ii p_aq_a)}{\sqrt{2\pi}},
 \qquad
 \braket{p,t}{q,t}
 =\prod_a\frac{\exp(-\ii p_aq_a)}{\sqrt{2\pi}},
\]
so
\[
 \braket{q',t+\dd t}{q,t}
 =\int\left(\prod_a\frac{\dd p_a}{2\pi}\right)
 \exp\left(
 \ii\sum_a p_a(q'_a-q_a)-H(q',p)\dd t
 \right).
\]
(The summation convention is suspended here.)

As in \hyperref[sec:9.1]{Section 9.1}, we divide the time interval from $-T/2$ to $T/2$ into a large number of very small intervals and insert a sum over $Q$ eigenstates for each interval. Defining functions $q(t)$ and $p(t)$ that interpolate between the values of the $Q$ and $P$ eigenvalues at each interval, we obtain the path integral expression for $F$ in its most general form
\begin{equation}
\begin{aligned}
 F(q',q;T)
 &=
 \int_{\substack{q(-T/2)=q\\q(T/2)=q'}}
 \left(\prod_{a,t}\dd q_a(t)\right)
 \int\left(\prod_{a,t}\frac{\dd p_a(t)}{2\pi}\right)\\
 &\quad{}\cdot
 \exp\left(
 \int_{-T/2}^{T/2}\dd t
 \left[
 \ii\sum_a\dot q_a(t)p_a(t)-H\bigl(q(t),p(t)\bigr)
 \right]
 \right).
\end{aligned}
 \tag{23.A.14}\label{eq:23.A.14}
\end{equation}
To calculate the partition function \eqref{eq:23.A.7} we would integrate over the $p$s and $q$s subject only to the condition that $q(t)$ is periodic with period equal to the inverse temperature $\beta$:
\begin{equation}
\begin{aligned}
 Z(\beta)
 &=
 \int_{q(\beta/2)=q(-\beta/2)}
 \left(\prod_{a,t}\dd q_a(t)\right)
 \int\left(\prod_{a,t}\frac{\dd p_a(t)}{2\pi}\right)\\
 &\quad{}\cdot
 \exp\left(
 \int_{-\beta/2}^{\beta/2}\dd t
 \left[
 \ii\sum_a\dot q_a(t)p_a(t)-H\bigl(q(t),p(t)\bigr)
 \right]
 \right).
\end{aligned}
 \tag{23.A.15}\label{eq:23.A.15}
\end{equation}
Eqs.~\eqref{eq:23.A.14} and \eqref{eq:23.A.15} look a little odd, with one term in the exponent real, and the other imaginary. These formulas begin to look more familiar after we do the path integral over the $p_a(t)$s. This integral is trivial in the important class of theories where $H(q,p)$ is quadratic in the $P$s:
\begin{equation}
 H(q,p)=\frac{1}{2}\sum_{a,b}A_{ab}(q)p_ap_b
 +\sum_aB_a(q)p_a+C(q).
 \tag{23.A.16}\label{eq:23.A.16}
\end{equation}
As shown in the \hyperref[app:9]{appendix to Chapter 9}, the integral over the $p$s in Eq.~\eqref{eq:23.A.14} yields
\begin{equation}
\begin{aligned}
 F(q',q;T)
 &=
 \int_{\substack{q(-T/2)=q\\q(T/2)=q'}}
 \left(\prod_{a,t}\dd q_a(t)\right)
 \left(\operatorname{Det}\left[2\ii\pi\mathcal A(q)\right]\right)^{-1/2}\\
 &\quad{}\cdot
 \exp\left(
 \int_{-T/2}^{T/2}\dd t
 \left[
 \ii\sum_a\dot q_a(t)\bar p_a(t)
 -H\bigl(q(t),\bar p(t)\bigr)
 \right]
 \right),
\end{aligned}
 \tag{23.A.17}\label{eq:23.A.17}
\end{equation}
where $\mathcal A(q)$ is the `matrix'
\begin{equation}
 \mathcal A(q)_{at',bt}\equiv\delta(t'-t)A_{ab}\bigl(q(t)\bigr),
 \tag{23.A.18}\label{eq:23.A.18}
\end{equation}
and $\bar p(t)$ is the stationary `point' of the argument of the exponential in Eq.~\eqref{eq:23.A.17}---that is, the solution of the equation
\begin{equation}
 \ii\dot q_a(t)
 =\left.
 \frac{\delta H(q(t),p)}{\delta p_a}
 \right|_{p=\bar p(t)}.
 \tag{23.A.19}\label{eq:23.A.19}
\end{equation}
The factor $i$ here should not be surprising, because Eq.~\eqref{eq:23.A.19} is the same equation that is satisfied by the non-Hermitian operators \eqref{eq:23.A.8}:
\[
 \ii\dot Q_a(t)=\ii[H,Q_a(t)]
 =\frac{\delta H\bigl(Q(t),P(t)\bigr)}{\delta P_a(t)}.
\]
For a Hamiltonian in the general quadratic form \eqref{eq:23.A.16}, the solution of Eq.~\eqref{eq:23.A.19} is
\begin{equation}
 \bar p_a=\sum_b\left[A^{-1}(q)\right]_{ab}
 \bigl(\ii\dot q_b-B_b(q)\bigr),
 \tag{23.A.20}\label{eq:23.A.20}
\end{equation}
and so Eq.~\eqref{eq:23.A.17} takes the form
\begin{equation}
\begin{aligned}
 F(q',q;T)
 &=
 \int_{\substack{q(0)=q\\q(T)=q'}}
 \left(\prod_{a,t}\dd q_a(t)\right)
 \left(\operatorname{Det}\left[2\ii\pi\mathcal A(q)\right]\right)^{-1/2}
 \exp\bigl(-S[q]\bigr),
\end{aligned}
 \tag{23.A.21}\label{eq:23.A.21}
\end{equation}
where $S[q]$ is the action
\begin{equation}
\begin{aligned}
 S[q]\equiv\int_{-T/2}^{T/2}\dd t\biggl[
 &\frac{1}{2}\sum_{a,b}A^{-1}_{ab}(q)\dot q_a\dot q_b
 +\ii\sum_{a,b}A^{-1}_{ab}(q)B_a(q)\dot q_b\\
 &-\frac{1}{2}\sum_{a,b}A^{-1}_{ab}(q)B_a(q)B_b(q)+C(q)
 \biggr].
\end{aligned}
 \tag{23.A.22}\label{eq:23.A.22}
\end{equation}
In the special (but common) case where $B_a(q)=0$, this simplifies to
\begin{equation}
 S[q]\equiv\int_{-T/2}^{T/2}\dd t
 \left[
 \frac{1}{2}\sum_{a,b}A^{-1}_{ab}(q)\dot q_a\dot q_b+C(q)
 \right].
 \tag{23.A.23}\label{eq:23.A.23}
\end{equation}
Thus in this case the `Lagrangian' appearing in the path integral is equal to what the Hamiltonian would be in Minkowskian spacetime when the $p$s are expressed in terms of $q$s and $\dot q$s.
